\begin{proposition}[扫描谱尾比的相位--深度同时提取：无需全反演即可定位主导缺陷]\label{prop:xi-scan-tail-ratio-dominant-defect}
\leanverified{paper\_xi\_scan\_tail\_ratio\_dominant\_defect}
在命题 \ref{con:xi-2kappa-minimal-complete-samples} 的指数谱表示下，设
\[
u_n=\sum_{j=1}^{\kappa}c_j\,r_j^{n}\qquad(n\ge 0),
\qquad
r_j=e^{-(1+\delta_j)\Delta}e^{-\mathrm{i}\gamma_j\Delta},
\qquad c_j>0.
\]
并假设模长严格分离：存在排序使得
\[
|r_1|>|r_2|\ge\cdots\ge|r_\kappa|.
\]
则存在极限
\[
\boxed{\;\lim_{n\to\infty}\frac{u_{n+1}}{u_n}=r_1\;}
\]
并且收敛为指数级：存在常数 \(C>0\) 使
\[
\left|\frac{u_{n+1}}{u_n}-r_1\right|
\ \le\ C\left(\frac{|r_2|}{|r_1|}\right)^{n}\qquad(n\to\infty).
\]
因此主导缺陷的深度与高度同余类可由尾部比值同时恢复：
\[
\boxed{\;
\delta_{\min}=\frac{-\log|r_1|}{\Delta}-1,
\qquad
\gamma_1\equiv-\frac{1}{\Delta}\arg(r_1)\ \Bigl(\bmod\ \frac{2\pi}{\Delta}\Bigr)\; }.
\]
\end{proposition}
\begin{proof}
写作 \(u_n=c_1r_1^n(1+\varepsilon_n)\)，其中
\(
\varepsilon_n:=\sum_{j\ge 2}\frac{c_j}{c_1}\bigl(\frac{r_j}{r_1}\bigr)^n
\)。
由 \(|r_j/r_1|\le |r_2/r_1|<1\) 得 \(\varepsilon_n=O((|r_2|/|r_1|)^n)\)，从而
\[
\frac{u_{n+1}}{u_n}
=r_1\,\frac{1+\varepsilon_{n+1}}{1+\varepsilon_n}
=r_1\Bigl(1+O\bigl((|r_2|/|r_1|)^n\bigr)\Bigr).
\]
最后两式由 \(|r_1|=e^{-(1+\delta_1)\Delta}\) 与 \(\arg(r_1)\equiv-\gamma_1\Delta\ (\bmod\ 2\pi)\) 立得。证毕。
\end{proof}

\begin{proposition}[严格模长分离下的逐次剥离定理：尾部极限诱导的谱分解在有限步精确闭合]\label{prop:xi-scan-successive-peeling-closure}
在命题 \ref{prop:xi-scan-tail-ratio-dominant-defect} 的假设下，令
\[
r_1:=\lim_{n\to\infty}\frac{u_{n+1}}{u_n},
\qquad
c_1:=\lim_{n\to\infty}\frac{u_{n}}{r_1^{n}}.
\]
并定义剥离序列
\[
u_n^{(1)}:=u_n-c_1r_1^n.
\]
则 \(u_n^{(1)}=\sum_{j=2}^{\kappa}c_jr_j^n\) 且仍满足同型模长分离。
递归进行 \(\kappa\) 次，得到有限步精确闭合：存在一族 \(\{(c_j,r_j)\}_{j=1}^{\kappa}\) 由尾部极限逐次唯一确定，使
\[
u_n=\sum_{j=1}^{\kappa}c_jr_j^n\qquad(\forall n\ge 0).
\]
因此，在严格模长分离的制度子类中，完整缺陷谱无需“全局一次性反演”，而可由尾部极限驱动的有限步谱剥离精确恢复；其不可避免的病态性仅在模长碰撞 \(|r_j|\to|r_k|\) 时发生，并与命题 \ref{con:xi-collision-illposedness-sep-blowup} 的“碰撞即回收”不适定性机制一致。
\end{proposition}
\begin{proof}
由模长分离，命题 \ref{prop:xi-scan-tail-ratio-dominant-defect} 给出 \(r_1\)。
再由 \(u_n=c_1r_1^n+O(|r_2|^n)\) 得 \(c_1=\lim_{n\to\infty}u_n/r_1^n\) 存在。
代入定义即得剥离恒等式 \(u_n^{(1)}=\sum_{j\ge 2}c_jr_j^n\)。
对 \(u^{(1)}\) 重复上述论证得到 \(r_2,c_2\)，迭代至 \(\kappa\) 步即得全谱恢复与唯一性。证毕。
\end{proof}
\leanverified{paper\_xi\_scan\_successive\_peeling\_closure}

\endinput
